\section{What connects one step to the next}
\label{sec:state}

\subsection{From the EFSM variables to the state record}
\label{sec:state-transition}

A machine diagram shows the permitted control flow. A shared record supplies the information
passed between steps. Routing and dispatch, described in later sections, both read this record.

A configuration of an extended state machine pairs a control location with variable values
(\cref{def:configuration}). An agent's runtime represents these separately. The application record
holds the variables; the runtime tracks the control location, or which computation runs next. The
\emph{state schema} declares each variable's name and type. Every step reads these variables and
proposes updates to some of them. In our agent, the constructor call
\demoref{StateGraph(AgentState)} constructs a graph with \demoref{AgentState} as its schema.

We describe a step with two operations. A node reads the current state $s_t$ and obtains an
external result $e_t$, then returns a partial update $u_t$. The runtime merges the update into the
record, key by key:
\[
  u_t \;=\; \mathrm{node}(s_t,\, e_t), \qquad
  s_{t+1} \;=\; \mathrm{merge}(s_t,\, u_t).
\]
Here $e_t$ denotes a result obtained during the node's execution, such as a model reply or tool
observation. The node's actual parameter is the state alone. A field omitted from the update keeps
its value in $s_{t+1}$.

\begin{figure}[htbp]
  \centering
  \begin{adjustbox}{width=\linewidth}% One transition: a node reads state and an external result, returns a partial update, and the
% runtime merges that update key by key. The dashed edge is the reducer's left argument, which
% comes from the current state and never from the node.
\begin{tikzpicture}[node distance=6mm and 5mm]
  \node[rnode] (s0) {state $s_t$};
  \node[mnode, right=of s0] (node) {node};
  \node[gnode, right=of node] (upd) {partial\\update $u_t$};
  \node[rnode, right=15mm of upd] (red) {reducer,\\one per key};
  \node[rnode, right=of red] (s1) {state $s_{t+1}$};
  \node[tnode, align=center, below=8mm of node] (ext) {model reply or\\tool observation $e_t$};
  \draw[flow] (s0) -- (node);
  \draw[flow] (node) -- (upd);
  \draw[flow] (upd) -- node[lab, below] {right\\argument} (red);
  \draw[flow] (red) -- (s1);
  \draw[flow] (ext) -- (node);
  \draw[back] (s0.north) to[out=35, in=145] node[lab, above] {left argument} (red.north);
\end{tikzpicture}
\end{adjustbox}
  \caption{One transition. The node proposes an update, and the runtime merges it. Each key has a
  reducer whose left argument is the stored value and whose right argument is the value returned
  by the node.}
  \label{fig:state-transition}
\end{figure}

A node returns only the keys it updates, because ``the Node does not need to return the whole State
schema - just an update''~\citep{langgraph-graph-api}. Our controller returns three of the loop's
ten keys: \demoref{messages}, \demoref{model\_calls} and \demoref{usage}. On the two paths that stop
the run before a model call, it returns only \demoref{status} (\cref{fig:controller-update}).

\begin{figure*}[tbp]
  \centering
  \begin{adjustbox}{width=\textwidth}% The controller's three return sites (plan_and_validation/graph.py, controller) against the
% loop's ten-key record. Each path returns a partial update; the runtime merges only the keys
% present, through that key's reducer, and every other key keeps its value.
\begin{tikzpicture}[
  key/.style={rnode, font=\ttfamily\scriptsize, minimum height=5mm, minimum width=13mm, inner sep=3pt},
  kept/.style={key, text=black!50, draw=ibmfg!35},
  hit/.style={key, draw=ibmblue, thick},
  upd/.style={gnode, font=\ttfamily\scriptsize, align=left, inner sep=4pt},
  flow/.style={-{Stealth[length=2mm]}, thick, draw=ibmfg!80, rounded corners=1.5mm},
]
  % ---- s_{t+1}: the record after the merge, ten keys in schema order ----------------------
  \node[kept] (k1) at (0,0) {alert};
  \node[kept, right=1mm of k1] (k2) {repo};
  \node[kept, right=1mm of k2] (k3) {commit};
  \node[kept, right=1mm of k3] (k4) {workspace};
  \node[hit,  right=1mm of k4] (k5) {messages};
  \node[kept, right=1mm of k5] (k6) {observations};
  \node[kept, right=1mm of k6] (k7) {candidate\_patch};
  \node[hit,  right=1mm of k7] (k8) {usage};
  \node[hit,  right=1mm of k8] (k9) {model\_calls};
  \node[hit,  right=1mm of k9] (k10) {status};
  \node[lab, font=\sffamily\footnotesize, left=3mm of k1] (s1) {state $s_{t+1}$};
  \node[lab, right=2mm of k10] {no reducer: replace};

  % ---- the three partial updates ----------------------------------------------------------
  \node[upd, anchor=south] (u3) at ($(k6.north)+(0,15mm)$)
    {\{"messages": [reply],\\\ "model\_calls": 1,\\\ "usage": \{input, output\}\}};
  \node[upd, anchor=south] (u2) at ($(k10.north)+(-2.5mm,15mm)$)
    {\{"status":\\\ NO\_PROGRESS\}};
  \node[upd, anchor=south west] (u1) at ($(u2.south east)+(3mm,0)$)
    {\{"status":\\\ BUDGET\_EXHAUSTED\}};

  % ---- the controller and the state it reads --------------------------------------------
  \node[mnode, anchor=south] (ctrl) at ($(u3.north)+(0,12mm)$) {controller};
  \node[rnode, left=16mm of ctrl] (s0) {state $s_t$};
  \draw[flow] (s0) -- (ctrl);

  % ---- controller to update: which path, and the guard that selects it --------------------
  \draw[flow] (ctrl.south) -- node[lab, right] {otherwise: one model call} (u3.north);
  \draw[flow] ([yshift=-1.5mm]ctrl.east) -| (u2.north)
    node[lab, pos=0.22, below] {last three observations identical};
  \draw[flow] ([yshift=1.5mm]ctrl.east) -| (u1.north)
    node[lab, pos=0.22, above] {model\_calls $\geq$ 40};

  % ---- update to record: one reducer per key that appears in the update -------------------
  \draw[flow] ([xshift=-10mm]u3.south) -- ++(0,-5mm) -| (k5.north)
    node[lab, pos=0.3] {add\_messages: append};
  \draw[flow] ([xshift=10mm]u3.south) -- ++(0,-9mm) -| (k8.north)
    node[lab, pos=0.3] {add\_usage: sum};
  \draw[flow] ([xshift=18mm]u3.south) -- ++(0,-4mm) -| (k9.north)
    node[lab, pos=0.3] {operator.add: $+1$};
  \draw[flow] (u2.south) -- ([xshift=-2.5mm]k10.north);
  \draw[flow] (u1.south) -- ++(0,-7mm) -| ([xshift=2.5mm]k10.north);

  % ---- the keys no path writes ------------------------------------------------------------
  \node[lab, text=black!60, anchor=north] at ($(k2.south)!0.5!(k3.south)+(0,-1.5mm)$)
    {kept from $s_t$: absent from every update};
  \node[lab, text=black!60, anchor=north] at ($(k6.south)!0.5!(k7.south)+(0,-1.5mm)$)
    {kept from $s_t$};
\end{tikzpicture}
\end{adjustbox}
  \caption{The controller's three return sites against the loop's ten-key record. The guard on
  each edge selects the path. The path that calls the model returns three keys, and each of the
  two paths that stop the run returns one. An arrow into a key names the reducer the runtime
  applies to it. \demoref{status} has no reducer, so the returned value replaces the stored one.
  The six keys without an arrow are absent from every update and keep their values from $s_t$.}
  \label{fig:controller-update}
\end{figure*}

The runtime applies the merge rule. It calls one reducer per updated key and saves the result. A
node can therefore return its contribution without implementing the choice between appending to a
list and replacing its stored value.

\subsection{Model context, runtime state, durable state}
\label{sec:state-three}

Model context and runtime state contain much of the same information. The model receives what it
needs for the current call. The runtime also retains counters and other fields that control
execution. Saving enough state to resume after the process ends makes that state durable.

\begin{table}[htbp]
  \centering
  \begin{adjustbox}{width=\linewidth}\begin{tabularx}{\linewidth}{@{}L{0.19\linewidth}X@{}}
    \toprule
    \textbf{Record} & \textbf{What it holds, and who reads it} \\
    \midrule
    Model context & Instructions, selected history and observations assembled for one model call.
    An agent with several model roles builds a separate context for each. In ours, the controller
    receives the conversation; the plan node and validator receive inputs assembled for their
    roles. \\
    Runtime state & The record retained between steps, including fields excluded from model calls.
    Nodes, routing functions and limit checks read it. \\
    Durable state & State saved beyond the process lifetime in a form that supports resuming a
    run. Resumption also requires the matching external state, such as the workspace, and a record
    of pending work. \\
    \bottomrule
  \end{tabularx}\end{adjustbox}
  \caption{Three kinds of record. Each call's model context is derived from runtime state.
  Durable state is runtime state saved beyond the process lifetime.}
  \label{tab:state-three}
\end{table}

Counters illustrate runtime information excluded from model context. Our loop refuses a model
call when \demoref{model\_calls} reaches forty. It also refuses a call when the last three
observations are identical. Neither number is shown to the model. Enforcement depends on the
runtime's comparison before the call, regardless of whether the model can see the limit in its
prompt.

If our agent's process stops during a run, a new process cannot recover its in-memory state
record. Both graphs are compiled without a checkpointer. Trace files and workspace artifacts may
survive, but they do not provide a resumable runtime record. LangGraph supports compiling a graph
with a checkpointer, which gives ``agents short-term
memory'' and ``lets LangGraph applications keep useful information beyond a single graph
run''~\citep{langgraph-persistence}. Our JSONL trace contains selected events
(\cref{sec:efsm-traces}). \textbf{The saved trace does not contain enough information to resume
the run.}

Here, \emph{durable} means that another process can read the saved record. Resuming can still
repeat work. What must be repeated depends on checkpoint and task boundaries; an interrupted node
may restart and repeat effects or paid calls not captured at a recoverable boundary. Later units
examine persistence and recovery. This lecture concerns the record needed to execute a run;
restart introduces further requirements.

\subsection{Reducers and the information retained between steps}
\label{sec:state-reducers}

A merge function determines whether a field accumulates a history or retains only its latest
value. This choice determines which evidence later steps can read.

In an extended state machine, each transition carries a set action that assigns new values to
variables (\citealp[slide~5]{seshia2015extended}). We implement that set action in two parts: a
step proposes an update, then the runtime combines it with the stored value.

\begin{definition}[Partial update]
\label{def:partial-update}
A partial update is the result of one step: a mapping from some of the state's fields to new
values. A field the update does not name keeps its stored value.
\end{definition}

\begin{definition}[Merge function]
\label{def:reducer}
A merge function belongs to one field. It takes the stored value and the partial update's value
as arguments, then returns the field's next value. Without a merge function, the update replaces
the stored value.
\end{definition}

LangGraph calls this merge function a \emph{reducer}, the term used throughout these notes. A
reducer has the form \demoref{(current, update) -> new} and is attached to a field with
\demoref{Annotated[T, fn]}. The runtime calls it whenever a node returns that key. LangGraph's
\demoref{add\_messages} reducer appends messages and replaces messages whose identifiers are
already stored. Other binary functions can also serve as reducers, including the standard
library's \demoref{operator.add} and functions we write ourselves.

Three merge behaviours cover most fields. Each determines what later steps can read.

\bi

\item \textbf{Append.} Add the update to the end of the stored list, keeping all earlier values.
This is needed when a later step examines a sequence, for example, to check whether the same
result occurred three times in a row. A field containing only the latest value cannot support
that check.

\item \textbf{Replace.} Overwrite the stored value with the update. This default behaviour suits a
current plan, decision or status, where the new value supersedes the previous one and determines
what later steps can use as evidence.

\item \textbf{Accumulate.} Combine the update arithmetically with the stored value, usually by
summing. Under this protocol, each node reports its own contribution, such as the cost of one
call. The reducer maintains the running total from those deltas.

\ei

The reducer also determines whether an update can clear a field. For a list that appends updates,
an empty update adds nothing: merging \demoref{[]} into \demoref{[x, y]} produces
\demoref{[x, y]}. \textbf{An empty update leaves an appended field unchanged.} Clearing it requires
bypassing the reducer and writing the new value directly, which LangGraph supports through
\demoref{Overwrite}. Such a reset should produce a separate trace event because it bypasses the
field's ordinary merge rule.

\begin{figure*}[tbp]
  \centering\footnotesize
  \begin{minipage}[t]{0.485\textwidth}
    \lfcaption{The loop}{reactive/state.py}
    \lfhead{class AgentState(TypedDict, total=False):}
    \lf{lifeid}{alert: dict}
    \lf{lifeid}{repo: str}
    \lf{lifeid}{commit: str}
    \lf{lifeid}{workspace: str}
    \lfgap
    \lf{liferun}{messages: Annotated[list, add\_messages]}
    \lf{liferun}{observations: Annotated[list, operator.add]}
    \lf{lifecur}{candidate\_patch: str | None}
    \lf{liferun}{usage: Annotated[Usage, add\_usage]}
    \lf{liferun}{model\_calls: Annotated[int, operator.add]}
    \lf{lifecur}{status: str}
  \end{minipage}\hfill
  \begin{minipage}[t]{0.485\textwidth}
    \lfcaption{The planning agent}{plan\_and\_validation/state.py}
    \lfhead{class AgentState(TypedDict, total=False):}
    \lf{lifeid}{alert: dict}
    \lf{lifeid}{repo: str}
    \lf{lifeid}{commit: str}
    \lf{lifeid}{workspace: str}
    \lfgap
    \lf[+]{lifecur}{plan: list | None}
    \lf[+]{lifehand}{failed\_step: dict | None}
    \lf[+]{lifewin}{replanned\_at: int}
    \lf{liferun}{messages: Annotated[list, add\_messages]}
    \lf{liferun}{observations: Annotated[list, operator.add]}
    \lf{lifecur}{candidate\_patch: str | None}
    \lf[+]{lifecur}{checks: dict | None}
    \lf[+]{lifecur}{validation: dict | None}
    \lf{liferun}{usage: Annotated[Usage, add\_usage]}
    \lf{liferun}{model\_calls: Annotated[int, operator.add]}
    \lf{lifecur}{status: str}
  \end{minipage}

  \medskip
  {\sffamily\scriptsize Lifetime:\enspace
  \lfkey{lifeid}{identity}\enspace\lfkey{liferun}{run}\enspace\lfkey{lifecur}{current}\enspace
  \lfkey{lifehand}{handoff}\enspace\lfkey{lifewin}{window}\qquad \textbf{+}\,new in the planning agent}
  \caption{The two \demoref{AgentState} classes, with comments omitted. Shading indicates field
  lifetime (\cref{sec:state-lifetimes}); \textbf{+} marks the planning agent's five additional
  fields. Three use the current lifetime already present in the loop. \demoref{failed\_step} and
  \demoref{replanned\_at} introduce two further lifetimes. Neither class has an attempt-lifetime
  field, because neither agent retries.}
  \label{fig:agentstate}
\end{figure*}

\subsection{Field lifetimes}
\label{sec:state-lifetimes}

A reducer defines how an update merges into a field. It does not define how long the resulting
value remains valid. A stale value can mislead later steps.

A field's \emph{lifetime} is the part of a run during which its value remains valid evidence for
its readers. It begins when the value is written and ends when an event supersedes or clears it.
We choose the lifetime when designing the field. The runtime enforces only the reducer, so the
nodes that write the field must follow its lifetime rules.

To choose a lifetime, identify who writes the field and which event makes its value invalid.
These answers distinguish six lifetimes.

\bi

\item \textbf{Identity.} Set once in the initial state and never updated by a node, so it needs no
reducer. Use it for the task, target and workspace that define what the run concerns and where it
may act. Updating an identity field changes the task during the run.

\item \textbf{Run.} Valid throughout the run and only grows. A history uses an append reducer; a
ledger uses an accumulating reducer. Use it when later steps or limits need the complete history,
such as the conversation, observations, or calls and tokens spent.

\item \textbf{Current.} Valid until its writer runs again, using the default replace rule. Use it
for the latest plan, candidate, check result or status, when readers need the current value and
the previous one is no longer valid evidence.

\item \textbf{Handoff.} Written by one node for one reader, which clears it after acting. Use it
when a step must explain why the next step was invoked, such as the reason for revising a plan.
If the reader leaves it set, the next routing decision responds to the same signal again.

\item \textbf{Window.} A span of run-lifetime history whose start is stored in an index field. Use
it when a check needs only the events since a given point, while other readers still need the
full history. Replace the index to move the window; retain the history.

\item \textbf{Attempt.} Valid for one attempt and cleared before a retry. Its reducer must allow
this reset: a replaced field can receive its empty value, while an appended field needs
\demoref{Overwrite}. Use it for results from a failed attempt that would mislead the next one,
such as the candidate and its check results.

\ei

Counter lifetimes determine whether limits hold. A limit on attempts or model calls requires a
run-lifetime counter that continues across attempts. \textbf{An attempt-lifetime counter resets
after each failure, so it cannot enforce a limit across attempts.}

\subsection{Invariants}
\label{sec:state-invariants}

A reducer or lifetime concerns one field. A property such as a spending limit constrains the
whole run. Proving it requires considering every transition the machine can take.

\begin{definition}[Invariant]
\label{def:invariant}
``An invariant for a state machine is a predicate, $P$, on states, such that whenever $P(q)$ is true
of a state, $q$, and $q \to r$ for some state, $r$, then $P(r)$ holds''~\citep[p.~5]{mit6042-lec16}.
\end{definition}

For an extended state machine, the states in this definition are configurations $(q, v)$, and the
arrow denotes one transition. To establish a property of every reachable configuration by
induction, we prove two statements: i) it holds initially, and ii) every transition preserves it.
The second proof obligation covers every configuration satisfying the predicate, including
unreachable ones. A property that holds in every reachable configuration may therefore need a
stronger predicate before we can prove preservation by every transition.

Consider the property \emph{model calls stay within the call budget}. Let $B$ be the budget, $c$
the recorded call counter, and $k$ the number of completed model calls. The runtime can compare
only $c$; $k$ describes the run's history and is not a program variable. We want $k \leq B$ and
prove the stronger predicate $P \equiv (c = k \,\wedge\, c \leq B)$. Initially, no call has been made
and $c = 0$, so $P$ holds. Transitions must then meet three conditions.

\bi

\item \textbf{Admission} (\cref{def:admission}). Every transition that issues a model call has the
guard $c < B$. Completing the call increases $k$ by one, keeping it at most $B$.

\item \textbf{Counting.} The update recording a completed call adds exactly one to $c$. An
unreported call leaves $c$ below $k$, allowing admission to exceed $B$. Reporting a call twice
exhausts the budget early.

\item \textbf{Preservation.} A transition without a call leaves $c$ unchanged. No transition,
including plan revision, decreases it. This is the run-lifetime requirement from
\cref{sec:state-lifetimes}. Resetting $c$ to zero preserves $c \leq B$ but breaks $c = k$, allowing
later calls to exceed $B$.

\ei

The argument covers completed calls only. An interrupted call may incur cost without being
counted, for example when an exception interrupts it. $P$ does not cover that cost.

An accumulating reducer on $c$ establishes part of the argument: merging adds the reported delta
to the stored count. Nodes must still check the guard before calling and report exactly one call
per completed call. \textbf{The reducer accumulates reported calls; each caller must report them
correctly.} In our agent, $c$ is \demoref{model\_calls} and $B$ is forty.

\subsection{State schemas of the loop and the planning agent}
\label{sec:state-demo}

Our state schema is the Python class \demoref{AgentState}, declared as a \demoref{TypedDict}
(\cref{fig:agentstate}). Each key names a record field, and its annotation specifies the value's
type. An annotation of \demoref{Annotated[T, fn]} also attaches the reducer \demoref{fn}. Setting
\demoref{total=False} marks every key as optional for type checking. This declaration accommodates
the partial updates that nodes return.
\demoref{TypedDict} supplies declarations for type checkers~\citep{pep589}; it does not check
returned values at run time.


The graph uses the class when \demoref{StateGraph(AgentState)} reads it during construction. It
creates one channel per key, using the annotated reducer or replacement when none is specified.
The remaining graph operations use those channels:

\bi

\item \demoref{add\_node(name, fn)} registers a node: a function that receives the whole state and
returns a partial update.

\item \demoref{add\_edge(a, b)} fixes a transition from node \demoref{a} to node \demoref{b}.

\item \demoref{add\_conditional\_edges(a, route, \{...\})} registers a routing function to run after
\demoref{a}. It reads the updated state and returns the next node's name without writing state.

\item \demoref{compile()} checks the wiring and returns a runnable graph.
\demoref{invoke(initial\_state(\ldots))} starts the run with a value for every key, so nodes do not
read missing keys.

\ei

\begin{lstlisting}[basicstyle=\ttfamily\scriptsize]
graph = StateGraph(AgentState)
graph.add_node("plan", self.plan)
graph.add_node("tools", self.tools)        # ... and three more
graph.add_edge(START, "plan")
graph.add_conditional_edges("tools", self.after_tools,
                            {"controller": "controller", "plan": "plan"})
graph.add_edge("validate", END)
app = graph.compile()
\end{lstlisting}

The loop's schema (\cref{fig:agentstate}, left) uses three lifetimes. Identity fields name the
scanner alert, fixture, pinned commit and workspace copy. Run-lifetime fields retain the
information used by stall and budget checks: \demoref{messages} and \demoref{observations} append,
while \demoref{model\_calls} and \demoref{usage} accumulate. Our \demoref{add\_usage} reducer sums
four token counters separately, without a combined total. Cache reads cost a fraction of fresh
input, and a combined total would hide an increase in cache reads alongside a decrease in fresh
input. In both agents, the reply adapter populates only input and output counters. Both cache
counters are zero in the planning agent's two recorded runs. The current-lifetime fields,
\demoref{candidate\_patch} and \demoref{status}, hold the latest patch and execution status.

The planning agent (\cref{fig:agentstate}, right) has five additional fields, each introduced with
the node that writes it. The plan node writes \demoref{plan}, the submit node writes
\demoref{checks}, and the validator writes \demoref{validation}. All three use the current
lifetime because a new value supersedes the previous one. The other two fields use different
lifetimes.

\bi

\item \textbf{Handoff.} When the tools node detects a stall, it writes the reason to
\demoref{failed\_step}. The routing function reads this field and selects the plan node. That node
reads the reason, revises the plan and clears \demoref{failed\_step} to \demoref{None}. If it left
the field set, the next pass through \demoref{tools} would return to the plan node again.

\item \textbf{Window.} The \demoref{observations} list appends and is never cleared. After a
revision, its last three entries are still the identical results that triggered revision. The
plan node records the current observation count in \demoref{replanned\_at}, and the stall check
then reads only \demoref{observations[replanned\_at:]}. This field also enforces the revision
budget: a nonzero value indicates that revision has occurred, so a second stall ends the run.

\ei

Each field in the planning agent's schema exists because a later computation reads what an earlier one
wrote (\cref{tab:state-flow}).

\begin{table}[htbp]
  \centering
  \begin{adjustbox}{width=\linewidth}\begin{tabularx}{\linewidth}{@{}L{0.30\linewidth}X@{}}
    \toprule
    \textbf{Fields} & \textbf{Why they persist and which steps read them} \\
    \midrule
    \demoref{alert}, \demoref{repo}, \demoref{commit}, \demoref{workspace} & Identify the task and
    its permitted workspace throughout the run. The workspace path names external state; the files
    are not in the record. \\
    \demoref{plan} & The plan node writes it and adds a rendered message to the controller's
    conversation. A revision reads the previous plan. \\
    \demoref{failed\_step}, \demoref{replanned\_at} & The tools node records a stall; routing sends
    the run to the plan node, which clears \demoref{failed\_step} and sets \demoref{replanned\_at}
    to the revised plan's observation boundary. Stall checks read from that index onward. \\
    \demoref{messages} & The opening instructions, plan messages, controller replies, tool replies
    and check feedback form the controller's conversation. This field excludes the replies of the
    plan node and validator. \\
    \demoref{observations} & The tools node appends one structured result per tool call; the stall
    checks and the validator read them. \\
    \demoref{candidate\_patch}, \demoref{checks} & The submit node stores the patch and the check
    results, whether passed or failed. Routing tests the checks, and validation reads both. \\
    \demoref{validation}, \demoref{status} & The validator stores its verdict, its reasons, and
    whether the model or runtime decided. \demoref{status} also records runtime stop reasons.
    Routing and the code that runs the graph and reports its outcome read this status. \\
    \demoref{model\_calls}, \demoref{usage} & Each of the three model nodes adds one call and that
    call's tokens. Admission reads the call count; only the run's closing summary reads the
    tokens. \\
    \bottomrule
  \end{tabularx}\end{adjustbox}
  \caption{The planning agent's fields, grouped by the node that writes them and the
  computations that read it.}
  \label{tab:state-flow}
\end{table}

Two code excerpts show how the call counter separates model context from runtime state
(\cref{sec:state-three}). The planning agent's controller receives the conversation as context.
Before every model call, admission compares the counter with the configured call budget.
\textbf{Admission reads the call budget. The controller's context does not contain it.} None of
the planning agent's prompts mentions the budget.

\noindent\begin{minipage}{\linewidth}
\begin{lstlisting}[basicstyle=\ttfamily\scriptsize]
# controller: the model receives the conversation
reply = self._call(model, state["messages"],
                   role="controller")

# admission: run before every model call
if state["model_calls"] >= self.config.budget:
    return {"status": TerminalStatus.BUDGET_EXHAUSTED}
\end{lstlisting}
\end{minipage}

We define a state graph in the following order.

\be

\item Name the nodes and the routing decisions between them.

\item List the values that nodes and routing functions must read. These values define the schema.
Omit fields that no computation reads.

\item For each field, identify its writer and the event that invalidates its value. Choose its
lifetime using \cref{sec:state-lifetimes}.

\item Choose the reducer from the lifetime: none for identity, append or accumulate for run, and
default replace for current, handoff and window. An attempt field needs a reducer that permits
clearing.

\item Write the initial state so that every key has a value before the first node runs.

\item Write each node to return only its own keys. Add each schema field with the node that
populates it.

\ee

\subsection{Reading}

\citet{langgraph-graph-api}, the implementation reference for state, schemas and
reducers; and \citet{mit6042-lec16}, \S2
(p.~5), on invariants and the Invariant Theorem.

For reference while writing HW1: \citet{pep589} on what a \demoref{TypedDict} does and does not
check at run time, and \citet{langgraph-persistence} on the checkpointers, threads and stores used
for durability in \cref{tab:state-three}.

\paragraph{Looking ahead.} A routing function reads the merged record, including the node's latest
update and previously stored values. It selects the next node or ends execution. The next section
examines conditional edges, feedback to the plan node or controller, and three interpretations of
a trace: recorded acceptance, a declared stop, and no recorded ending.
